\section{METR 的失败模式与有效性边界}

\subsection{长轨迹为何失败}

\videofigure{figures/metr-failure-modes.jpg}{常见失败包括规划或工具选择不当、错误推理、提前放弃以及重复动作循环。}{00:23:18--00:26:18}

错误通常不是最后一步突然出现，而是从早期决策累积：选择了错误工具、没有检查中间产物、误读环境状态，随后不断在错误分支上投入 token 和时间。

\begin{warningbox}{更多 context 不能替代状态管理}
把完整历史塞回模型不等于 agent 记住了关键状态。长轨迹需要显式任务图、已验证事实、未解决问题、工具输出摘要和回滚点。
\end{warningbox}

\begin{knowledgebox}{循环行为是停止策略失败}
Agent 反复执行同一动作，往往说明它没有把“动作已尝试但无效”写入状态，也没有触发升级、改计划或终止。检测重复状态比单纯增加最大步数更有效。
\end{knowledgebox}

\subsection{任务是否真的代表现实工作}

\videofigure{figures/metr-validation.jpg}{METR 从环境杂乱度、任务难度偏差和内部评审等角度验证 time-horizon 结论。}{00:26:18--00:28:06}

验证需要检查：任务是否过于干净、自动评分是否漏掉质量问题、人类时间估计是否可靠、模型是否见过题目、环境错误是否被误算为能力错误。课程强调，评估可信度来自这些审计，而不是一条漂亮趋势线。

\begin{importantbox}{Benchmark 也需要 verifier}
Agent 被 verifier 评估，benchmark 本身同样要被验证：任务代表性、评分规则、重复运行方差和人类基线都可能成为系统性误差源。
\end{importantbox}

\subsection{本章小结}

METR 适合回答能力跨度，但不能独自回答工作产品质量、经济价值和现实上下文问题。分析失败轨迹与 benchmark 有效性，是正确解读 time horizon 的必要部分。

